% !TEX program = xelatex
\documentclass[a4paper, 11pt]{ctexart}
\usepackage[margin=2cm]{geometry}
\usepackage{xcolor}
\usepackage{booktabs}
\usepackage{titlesec}
\usepackage{tcolorbox}
\usepackage{enumitem}
\usepackage{hyperref}
\usepackage{float}
\usepackage{array}
\usepackage{amsmath}

\definecolor{tsinghua}{RGB}{102, 8, 116}
\definecolor{bglight}{RGB}{247, 245, 248}
\titleformat{\section}[block]{\Large\bfseries\color{tsinghua}}{\thesection}{1em}{}
\titleformat{\subsection}[block]{\large\bfseries\color{tsinghua!80!black}}{\thesubsection}{1em}{}
\linespread{1.35}
\setlist[itemize]{itemsep=4pt, parsep=0pt, topsep=4pt}
\hypersetup{colorlinks=true, linkcolor=tsinghua, urlcolor=tsinghua}

\title{\vspace{-1.5cm}\color{tsinghua}\textbf{方向 C：面向动力装备的专业知识库与智能体构建}\\\Large{—— "鸿雁计划"科创基金细化企划书}}
\author{\textbf{汇报人：纪文龙} \\ \small 清华大学行健书院 2024级本科生}
\date{\small 2026年4月}

\begin{document}
\maketitle
\thispagestyle{empty}

\begin{tcolorbox}[colback=tsinghua!5, colframe=tsinghua!50, title=\textbf{一句话概括}, fonttitle=\bfseries]
把动力设备领域几十年积累的专家经验"灌"进大模型，建一个\textbf{能查资料、能推理、能给维修建议}的AI助手——一线工程师用自然语言问它"这台机器怎么了？"就能得到靠谱答案。发表 1 篇 SCI Q1/Q2 论文。
\end{tcolorbox}

% ================================================================
\section{这个研究到底在干什么？（白话版）}
% ================================================================

\subsection{先理解"知识库"是什么}

想象一个老师傅在发电厂工作了30年，他脑子里装着海量的经验：哪种振动波形对应轴承磨损、什么温度曲线意味着快要出事、遇到这种故障应该先查哪个部件……这些经验就是\textbf{领域知识}。

问题是：老师傅会退休、会离职，经验带走了就没了。而且一个人的经验再丰富也有盲区。\textbf{知识库}就是把这些分散在技术文档、维修记录、行业标准、学术论文里的知识\textbf{数字化、结构化}，存进计算机里，让任何人都能检索和使用。

\subsection{那"智能体"是什么？}

你可能已经用过ChatGPT——它能聊天、能写文章，但它\textbf{不懂动力设备}。你问它"这台汽轮机的2X频振动偏高可能是什么原因？"它可能瞎编一个看似合理但实际错误的答案。这叫\textbf{幻觉（hallucination）}。

\textbf{智能体（Agent）} = 大模型 + 工具调用 + 知识检索。它不是光靠"脑子里"的知识回答，而是\textbf{先去知识库里找相关资料，再结合资料进行推理}。就好比一个医生不光靠记忆诊断，还要查检验报告、翻教科书——这样的诊断才靠谱。

\subsection{这两者结合起来能干什么？}

\begin{itemize}
    \item \textbf{故障诊断辅助}：一线工程师描述现象（"3号机组振动突然升高"），智能体检索类似历史案例、查阅设备手册，给出可能的故障原因和排查建议。
    \item \textbf{知识问答}：新入职的工程师不用翻几百页的技术手册，直接问智能体。
    \item \textbf{经验沉淀}：老专家的诊断记录、维修日志不再是"死"的文档，而是变成智能体可以检索和推理的"活"知识。
\end{itemize}

\subsection{所以我们要做的核心事情是：}

\begin{tcolorbox}[colback=bglight, colframe=tsinghua!30, arc=2mm]
\textbf{针对动力设备（汽轮机、发电机、轴承等）场景，（1）构建一个专业知识库——把技术文档、故障案例、运行规律整理成大模型能"读懂"的格式；（2）在知识库上搭建一个RAG智能体——让它在回答问题时先检索知识库再生成答案，消除幻觉、保证专业性。}
\end{tcolorbox}


\section{为什么现在做这件事正当时？}

\subsection{大模型能力已经成熟，但工业落地还有"最后一公里"}

2023--2025年，GPT-4、Claude、Qwen等大模型的通用能力已经很强了。但工业界的反馈是：\textbf{通用大模型直接用到专业领域不靠谱}——回答不够精准、容易编造、缺乏可追溯性。

\subsection{RAG（检索增强生成）解决了"幻觉"问题}

RAG的核心思想很简单：模型回答问题之前，先从一个可信的知识库中\textbf{检索出相关段落}，把这些段落作为"参考资料"附在提示词里，然后让模型\textbf{基于这些资料}来生成答案。这样模型的回答就有据可查，不再是"空口说白话"。

\textbf{技术上包括三步}：
\begin{enumerate}
    \item \textbf{索引（Indexing）}：把文档切成小段（chunk），用Embedding模型把每段文字变成一个向量（一串数字），存进向量数据库。
    \item \textbf{检索（Retrieval）}：用户提问也转成向量，找到知识库中语义最接近的几段。
    \item \textbf{生成（Generation）}：把检索到的段落和用户问题一起交给大模型，生成最终答案。
\end{enumerate}

\subsection{从"基础RAG"到"Agentic RAG"——智能体化}

基础RAG是"被动的"：用户问什么就搜什么。但真实的故障诊断需要\textbf{多步推理}：先确定症状 → 缩小可能原因 → 逐步排查 → 给出建议。这就需要一个\textbf{Agent框架}让大模型能够"思考、规划、调用工具"。

\begin{tcolorbox}[colback=bglight, colframe=tsinghua!30, arc=2mm]
\textbf{当前文献空白}：大模型+RAG的工业应用论文很多，但\textbf{针对能源动力设备}（汽轮机、发电机组等）的系统性知识库和智能体工作还很少。刘老师课题组在故障诊断领域积累了大量数据和方法论，正好可以填补这个空白。
\end{tcolorbox}


\section{我们的方法：PowerKB-Agent}

\subsection{整体架构（用大白话讲）}

我们的系统叫 \textbf{PowerKB-Agent}（Power equipment Knowledge Base Agent），分三层：

\begin{enumerate}
    \item \textbf{知识层（Knowledge Layer）}：
    \begin{itemize}
        \item 收集动力设备的技术文档（设备手册、国标、行业规范）、故障案例库（历史维修记录、典型故障分析报告）、学术论文（刘老师课题组发表的高被引论文等）。
        \item 对文档进行\textbf{清洗-分块-标注}：清洗掉无关内容（页眉页脚、广告等），按语义切成合理大小的文本块，给每个块打上元数据标签（设备类型、故障类别、知识类型等）。
        \item 用Embedding模型（如BGE-large-zh）将文本块编码为向量，存入向量数据库（Milvus或Chroma）。
    \end{itemize}

    \item \textbf{检索层（Retrieval Layer）}：
    \begin{itemize}
        \item 用户提问 → 转成向量 → 在向量数据库中找Top-K最相关的文本块。
        \item 加入\textbf{重排序器（Reranker）}：粗检索找回20个候选，再用一个小模型精排，选出最相关的5个——提高检索精度。
        \item 支持\textbf{混合检索}：向量相似度 + 关键词匹配（BM25），两种方式互补，确保既能找到语义相近的，也能找到术语完全匹配的。
    \end{itemize}

    \item \textbf{智能体层（Agent Layer）}：
    \begin{itemize}
        \item 基于LangChain/LlamaIndex搭建Agent框架。
        \item Agent能做的事：（a）检索知识库；（b）调用故障诊断规则（if-then专家规则）；（c）查询结构化数据库（设备参数表、历史运行数据）；（d）多轮对话追问用户以获取更多信息。
        \item 最终输出带\textbf{来源引用}的回答——每个诊断建议都标注"这条信息来自XX文档第XX页"，可追溯、可验证。
    \end{itemize}
\end{enumerate}

\subsection{为什么这样设计有效？}

\begin{itemize}
    \item \textbf{知识库保证专业性}：答案基于真实的技术文档和案例，不是大模型"编"的。
    \item \textbf{RAG消除幻觉}：有检索做支撑，大模型不会胡说。
    \item \textbf{Agent实现多步推理}：不是简单的一问一答，而是能分步骤思考和排查。
    \item \textbf{引用溯源增强信任}：工程师可以直接点开引用来源核实，符合工业场景的可靠性要求。
\end{itemize}

\subsection{三层架构与工作量}
\begin{enumerate}
    \item \textbf{底层·数据}：课题组积累的动力设备故障案例库、技术文档、公开的国标/行标、刘老师团队发表的论文——这是我们相对于其他团队的\textbf{独家优势}。
    \item \textbf{中层·系统}：RAG Pipeline + Agent框架——这是论文的核心贡献，重点在知识库构建方法论和检索增强策略。
    \item \textbf{上层·展示}：一个对话式Web界面，输入故障描述就能得到诊断建议和引用来源。
\end{enumerate}


\section{实验设计（怎么证明系统有效？）}

\subsection{评估思路}

和方向A/B不同，这个方向不是"跑神经网络算准确率"，而是评估一个\textbf{问答系统}的质量。我们需要回答两个问题：
\begin{itemize}
    \item \textbf{检索质量}：给一个问题，系统找到的文档段落是否确实包含正确答案？
    \item \textbf{生成质量}：基于检索到的段落，模型最终给出的回答是否准确、完整、有用？
\end{itemize}

\subsection{评估数据集构建}

我们需要人工构建一个\textbf{测试题库}（约100--200道题），来源包括：
\begin{itemize}
    \item 从课题组的故障案例中提取真实问答对（"问：某机组出现XX现象，可能原因？答：……"）
    \item 从课题组论文中提取技术性问题（"问：DANN中梯度反转层的作用是什么？"）
    \item 从设备手册中提取操作性问题（"问：XX型轴承的允许温升范围？"）
\end{itemize}

\subsection{评价指标}

\begin{table}[H]
\centering
\renewcommand{\arraystretch}{1.2}
\small
\begin{tabular}{p{3cm} p{4cm} p{5.5cm}}
    \toprule
    \textbf{指标} & \textbf{含义} & \textbf{通俗解释} \\
    \midrule
    Hit Rate@K & 检索回来的Top-K段落中是否包含正确答案 & "前5个搜索结果里有没有答案" \\
    MRR & 正确答案在检索结果中的平均排名位置 & "答案一般排第几？越靠前越好" \\
    Faithfulness & 生成的回答是否与引用的段落内容一致 & "模型有没有偏离引用的资料" \\
    Answer Relevancy & 回答对问题的相关程度 & "有没有答非所问" \\
    专家评分 & 邀请刘老师/课题组博士生评分 & "行内人觉得这个回答靠谱吗" \\
    \bottomrule
\end{tabular}
\end{table}

\subsection{消融实验}

逐步去掉系统的组件，看性能变化：
\begin{itemize}
    \item 去掉RAG（= 纯大模型直接回答） $\to$ 证明知识库检索的价值
    \item 去掉Reranker（= 只用粗检索） $\to$ 证明重排序的价值
    \item 去掉混合检索（= 只用向量检索） $\to$ 证明关键词匹配的互补作用
    \item 去掉Agent（= 单轮RAG） $\to$ 证明多步推理的价值
    \item 去掉元数据过滤（= 不按设备/故障类型筛选） $\to$ 证明结构化标注的价值
\end{itemize}

\subsection{对比基线}

\begin{itemize}
    \item \textbf{纯大模型}：直接用GPT-4/Qwen回答专业问题（不检索）
    \item \textbf{基础RAG}：Naive RAG（固定大小分块 + 纯向量检索 + 单轮生成）
    \item \textbf{Google搜索}：用搜索引擎搜同样的问题，看返回结果质量
    \item \textbf{传统专家系统}：基于if-then规则的故障诊断系统
\end{itemize}

只要PowerKB-Agent在Faithfulness和专家评分上明显优于上述基线，论文就有说服力。


\section{预期成果}
\begin{itemize}
    \item \textbf{1篇SCI Q1/Q2论文}：投向 \textit{Expert Systems with Applications}（IF$\sim$8.5）或 \textit{Knowledge-Based Systems}（IF$\sim$8.1）。"领域知识库+RAG+Agent应用于能源动力设备智能运维"是当前的热点空白。
    \item \textbf{动力装备知识库}：一个结构化、可检索的专业知识库（含文档$>$500篇、标注问答对$>$200组），可供课题组长期使用和迭代。
    \item \textbf{智能体演示系统}：一个Web界面的对话式诊断助手，输入故障描述即可获得带引用溯源的诊断建议。
    \item \textbf{鸿雁计划结题}：完整结题报告 + 上述成果。
\end{itemize}


\section{W1--W30 逐周工作计划}

每周投入 \textbf{8 小时}，分 6 阶段。

\subsection{阶段一：调研 + 认识领域（W1--W5，4月）}

\noindent\textit{这个阶段的目标是}：搞懂大模型+RAG+Agent在工业领域的应用现状，同时摸清课题组有哪些可用的知识资料。

\begin{description}[style=nextline, leftmargin=1.2cm, labelwidth=1cm, font=\bfseries\color{tsinghua}]
\item[W1] \textbf{读20篇相关论文}。搜索"LLM industrial fault diagnosis""RAG knowledge base power equipment""Agentic RAG maintenance"，做方法对比表（架构/知识来源/评估方式/效果/年份）。重点关注：别人的知识库怎么建的？RAG的检索策略用了什么？Agent的工具调用设计了哪些？
\item[W2] \textbf{精读3--5篇核心论文}。(1) RAG综述（最新一篇Survey）；(2) 工业领域Agent应用论文；(3) 刘老师发表的故障诊断方法论文（理解课题组的核心技术体系）。输出：一份2页的技术选型报告发刘老师确认（用哪个大模型、哪个向量数据库、哪个Agent框架）。
\item[W3] \textbf{盘点课题组知识资源}。和课题组博士生沟通，梳理可用的文档资料：设备技术手册有哪些？历史故障案例库存在什么地方？格式是什么（PDF/Word/数据库记录）？整理一份"知识资源清单"。
\item[W4] \textbf{学习工具链}。动手跑通一个最简单的RAG Demo：用LangChain加载几份PDF文档 → 切分 → 存入Chroma向量数据库 → 问几个问题看效果。目的是把整个Pipeline的每一步都亲手走一遍。
\item[W5] \textbf{学习Agent框架}。在W4的基础上给Agent加"工具"：比如一个能查设备参数表的工具、一个能搜索故障案例的工具。确认自己能用LangChain/LlamaIndex搭出一个带工具调用的Agent。
\end{description}

\subsection{阶段二：知识库构建（W6--W10，5--6月）}

\noindent\textit{这个阶段的目标是}：把课题组的知识资料真正地"灌"进系统里，建成一个可检索的专业知识库。

\begin{description}[style=nextline, leftmargin=1.2cm, labelwidth=1cm, font=\bfseries\color{tsinghua}]
\item[W6] \textbf{文档预处理Pipeline}。开发一套自动化脚本：(1) PDF/Word → 纯文本（用PyMuPDF或Unstructured库处理排版、表格、公式）；(2) 文本清洗（去页眉页脚、修复乱码）；(3) 智能分块（按段落/章节语义切分，而非固定字数切），每块附加元数据（文档标题、章节名、设备类型、日期）。先处理50份文档，验证Pipeline可用性。
\item[W7] \textbf{大规模文档编码入库}。用BGE-large-zh（或其他中文Embedding模型）对所有文本块生成向量，存入向量数据库。同时建BM25倒排索引用于关键词检索。目标：处理完全部可用文档（预计300--500份）。
\item[W8] \textbf{知识标注与质量检查}。(1) 对重要文档块人工标注"设备类型-故障模式-知识类型"三级标签，支持后续按类别过滤检索；(2) 从知识库中随机抽50个问题做初步检索测试，看Top-5命中率。
\item[W9] \textbf{构建评估题库}。从故障案例和论文中人工编写100--200道测试问题（含标准答案和引用来源），包括：事实型（"XX设备的额定参数是？"）、诊断型（"出现XX现象可能是什么故障？"）、方法型（"如何用DANN做跨域迁移？"）三类。
\item[W10] \textbf{阶段验证 + 汇报}。用测试题库评估"基础RAG"的效果——Hit Rate@5、MRR等指标。把结果发刘老师，确认知识库基础质量合格后进入下一阶段。
\end{description}

\subsection{阶段三：核心研发——RAG优化与Agent搭建（W11--W15，7月）}

\noindent\textit{这是最关键的阶段}：从"能跑"提升到"跑得好"，实现论文的核心创新。

\begin{description}[style=nextline, leftmargin=1.2cm, labelwidth=1cm, font=\bfseries\color{tsinghua}]
\item[W11] \textbf{检索优化}。(1) 实现混合检索（向量+BM25加权融合）；(2) 加入Reranker模型（bge-reranker-large）做二次精排；(3) 实现Query Expansion——把用户短问题用LLM扩展成更详细的检索Query（比如"振动高"→"旋转机械振动幅值偏高可能原因排查"），提升召回率。
\item[W12] \textbf{搭建Agent框架}。定义Agent可调用的工具集：(1) 知识库向量检索工具；(2) 设备参数结构化查询工具（SQL查设备型号/参数表）；(3) 故障规则推理工具（基于课题组总结的专家规则）；(4) 追问工具（当信息不足时主动向用户追问）。用LangGraph或AutoGen实现Agent的"规划-执行-反思"循环。
\item[W13] \textbf{Prompt工程}。精心设计System Prompt：让智能体扮演"动力设备高级诊断工程师"角色，明确回答格式要求（必须列出故障可能性、排查建议、引用来源），设置安全规则（不确定时明确说"我不确定"而非编造）。测试不同Prompt对回答质量的影响。
\item[W14] \textbf{端到端联调}。把知识层、检索层、Agent层串联起来，用评估题库跑完整的端到端测试。记录全部指标，找出短板（是检索不准？还是生成不好？还是Agent逻辑有问题？）。
\item[W15] \textbf{调优 + 向刘老师汇报}。针对短板逐一优化（调分块大小、检索Top-K、Reranker阈值、Prompt措辞等），写中期报告。
\end{description}

\subsection{阶段四：大规模实验——出论文数据（W16--W20，8--9月）}

\noindent\textit{这个阶段的目标是}：跑完所有对比实验和消融实验，产出论文所需的全部数据和图表。

\begin{description}[style=nextline, leftmargin=1.2cm, labelwidth=1cm, font=\bfseries\color{tsinghua}]
\item[W16] \textbf{消融实验}。5种消融配置（纯LLM / 无Reranker / 无混合检索 / 无Agent / 无元数据过滤） × 完整评估题库，记录所有指标对比。
\item[W17] \textbf{基线对比实验}。跑4种基线（纯GPT-4、Naive RAG、Google搜索、传统专家系统）在同一题库上的表现，做完整对比表。
\item[W18] \textbf{专家评估}。邀请刘老师和2--3位课题组博士生对50道诊断型问题的回答做盲评打分（1--5分的准确性、完整性、实用性），计算评分间一致性（Cohen's Kappa）。
\item[W19] \textbf{案例分析}。选5个典型故障诊断场景，展示Agent的完整推理过程（用户提问 → Agent规划检索 → 获取信息 → 逐步推理 → 给出建议），做成论文中的Case Study图。
\item[W20] \textbf{结果整理}。汇总全部实验数据，画论文图表：性能对比柱状图、消融热力图、Agent推理流程图、知识库统计信息表等。
\end{description}

\subsection{阶段五：写论文 + 搭演示系统（W21--W25，10--11月）}

\begin{description}[style=nextline, leftmargin=1.2cm, labelwidth=1cm, font=\bfseries\color{tsinghua}]
\item[W21] \textbf{搭建Web演示界面}。用Streamlit或Gradio做一个对话式界面：左侧输入故障描述、右侧显示智能体回答（含引用来源链接），底部显示Agent的"思考过程"（检索了什么、调用了什么工具）。
\item[W22] 和刘老师确认论文结构和目标期刊；写论文第1章（动力设备智能运维的背景和需求、为什么需要领域知识库+智能体）。
\item[W23] 写论文第2章（相关工作：RAG技术综述、工业领域LLM应用现状、知识库构建方法论）。
\item[W24] 写论文第3章（PowerKB-Agent系统设计：知识库构建流程、检索增强策略、Agent架构，含架构图和算法伪代码）。
\item[W25] 写论文第4章（实验结果：全部对比表、消融分析、案例分析、专家评估结果）。
\end{description}

\subsection{阶段六：投稿 + 结题（W26--W30，12--1月）}

\begin{description}[style=nextline, leftmargin=1.2cm, labelwidth=1cm, font=\bfseries\color{tsinghua}]
\item[W26] 第5章（结论、局限性讨论、未来展望）+ 全文润色。
\item[W27] 给刘老师审改，根据反馈修改补充。
\item[W28] 按目标期刊模板排版，提交投稿。
\item[W29] 写鸿雁计划结题报告（约15页）。
\item[W30] 提交结题材料包，项目完成。
\end{description}

\end{document}
